\documentclass[11pt,a4paper]{article}

\usepackage[top=24mm,bottom=23mm,left=27mm,right=27mm,headheight=22pt]{geometry}
\usepackage{microtype}
\usepackage{enumitem}
\usepackage{xcolor}
\usepackage{hyperref}
\usepackage{fancyhdr}

\definecolor{ink}{HTML}{1C2530}
\definecolor{softink}{HTML}{626B75}
\definecolor{accent}{HTML}{285FAE}
\definecolor{rulecolor}{HTML}{C9C3B8}

\hypersetup{
  pdftitle={Opening Note},
  pdfauthor={Xayah Hina},
  pdfsubject={An introduction to the Writing space at i.xayah.me},
  colorlinks=true,
  linkcolor=accent,
  urlcolor=accent
}

\setlength{\parindent}{0pt}
\setlength{\parskip}{0.62em}
\setlist{leftmargin=1.5em,itemsep=0.4em,topsep=0.5em}
\setcounter{secnumdepth}{0}

\pagestyle{fancy}
\fancyhf{}
\fancyhead[L]{\footnotesize\color{softink}\textsc{Xayah Hina / Writing}}
\fancyhead[R]{\footnotesize\color{softink}Opening Note}
\fancyfoot[C]{\footnotesize\color{softink}\thepage}
\renewcommand{\headrulewidth}{0.4pt}
\renewcommand{\headrule}{\hbox to\headwidth{\color{rulecolor}\leaders\hrule height \headrulewidth\hfill}}

\title{\textbf{Opening Note}}
\author{Xayah Hina}
\date{15 July 2026}

\begin{document}

\hypersetup{pageanchor=false}
\begin{titlepage}
  \thispagestyle{empty}

  {\small\bfseries\color{accent}\MakeUppercase{Writing No. 001}\par}

  \vspace*{0.2\textheight}

  {\fontsize{39}{43}\selectfont\bfseries\color{ink}Opening Note\par}

  \vspace{2.2em}
  {\color{rulecolor}\rule{\textwidth}{0.7pt}\par}
  \vspace{1.4em}

  {\large\color{softink}
    A beginning for a place about graphics, music, technical work,
    and the texture of ordinary days.\par
  }

  \vfill

  {\large\bfseries Xayah Hina\par}
  \vspace{0.35em}
  {\color{softink}15 July 2026\par}
  \vspace{0.35em}
  {\small\href{https://i.xayah.me/}{i.xayah.me}\par}
\end{titlepage}
\hypersetup{pageanchor=true}

\section{Why This Space Exists}

Some thoughts arrive in the middle of technical work and disappear before
they have a place to settle. Others begin with a piece of music, a visual
detail, or an ordinary moment that seems too small for a formal project and
too important to leave unrecorded. This site is a home for both kinds.

My academic page is meant to stay concise: research interests, publications,
projects, and the work that can be named clearly. This space can move at a
different pace. It can hold unfinished questions, explanations written after
the code is done, and observations whose value may only become clear later.

\section{What I Hope to Write}

The subjects will change, but they will probably return to a few recurring
threads:

\begin{itemize}
  \item \textbf{Graphics and research.} Notes on simulation, rendering,
    inverse design, tools, failed approaches, and the reasoning that rarely
    fits into a project page.
  \item \textbf{Music and listening.} Records of what I am making, hearing,
    and learning from sound.
  \item \textbf{Everyday life.} Small essays, travel notes, photographs in
    words, and moments worth keeping without forcing them into a larger claim.
\end{itemize}

\section{On Format}

Each entry is written in \LaTeX{} and published as a PDF. That choice is
intentionally a little slower than posting into a feed. Typesetting asks me to
decide where a thought begins and ends, while the PDF gives each entry a stable
form that can be read away from the site.

The web page remains an index rather than becoming another publishing system.
It tells you what is here; the documents carry the writing itself.

\section{A Beginning}

This first note is deliberately simple. It marks a beginning and leaves enough
empty space for the character of the site to emerge through future writing. If
something here resonates with you, you can find my academic work at
\href{https://graphics.xayah.me/}{graphics.xayah.me}.

\vfill
{\color{rulecolor}\rule{\textwidth}{0.5pt}\par}
{\small\color{softink}Typeset with Tectonic. Source and future notes live at
\href{https://i.xayah.me/}{i.xayah.me}.\par}

\end{document}
